\subsection{Alternate Multilinear Forms}

\begin{tcolorbox}[title=Problem 1, breakable]
    Let $f \in \mathcal{A}(V)$ then $f(x_1, \ldots, x_n) = 0$
    whenever $x_1, \ldots, x_n$ are linearly dependent.
\end{tcolorbox}

\begin{proof}
    Suppose $x_1, \ldots, x_n$ are linearly dependent.
    There exists scalars $c_1, \ldots, c_n$, not all zero, such that 
    $\sum c_i x_i = 0$. Suppose WLOG $c_j \ne 0$ then 
    \[x_j = -\sum_{i \ne j} \frac{c_i}{c_j} x_i.\]
    Then 
    \[
    f(x_1, \ldots, x_j, \ldots, x_n)
    =
    f\left(x_1, \ldots, -\sum_{i \ne j} \frac{c_i}{c_j} x_i, \ldots, x_n\right).
    \]
    By multilinearity,
    \[
    f(x_1, \ldots, x_n)
    =
    -\sum_{i \ne j} \frac{c_i}{c_j}
    f(x_1, \ldots, x_i, \ldots, x_n),
    \]
    where in each term $x_i$ appears in both the $i$th and $j$th positions.
    Since $f$ is alternating, each such term is $0$. Thus
    \[
    f(x_1, \ldots, x_n)=0.
    \]
\end{proof}

\newpage
\begin{tcolorbox}[title=Problem 2, breakable]
\begin{enumerate}[(i)]
    \item If $f_1, \ldots, f_n$ are linear forms on $V$, then the function $f: V^n \to F$ defined by
    \[
        f(x_1, \ldots, x_n) = f_1(x_1) \cdots f_n(x_n)
    \]
    is multilinear.

    \item If $n = 2$ and $f_1, f_2$ are nonzero linear forms on $V$, then the bilinear form $f$ of (i) is not alternate. 
    \{Hint: \S1.6, Exercise 18.\}

    \item If $n = 2$ and $f_1, f_2$ are linearly independent linear forms on $V$, then the bilinear form $f$ defined by
    \[
        f(x_1, x_2) = f_1(x_1)f_2(x_2) - f_1(x_2)f_2(x_1)
    \]
    is alternate and not identically zero. \{Hint: $\ker f_1$ is not contained in $\ker f_2$.\}

    \item If $x_1, \ldots, x_n$ is a basis of $V$, and $f_1, \ldots, f_n$ is the dual basis of $V'$, then the function $f$ defined by the formula
    \[
        (*) \qquad f(y_1, \ldots, y_n) = \sum_{\sigma \in S_n} (\operatorname{sgn} \sigma) f_{\sigma(1)}(y_1) \cdots f_{\sigma(n)}(y_n)
    \]
    is the nonzero alternate multilinear form constructed in Lemma 7.1.13.
\end{enumerate}
\end{tcolorbox}

\begin{proof}
    Suppose $f_1, \ldots, f_n$ are linear forms on $V$.
    Multilinearity follows from the linearity of each $f_i$.
\end{proof}

\begin{proof}
    Suppose $n = 2$ and $f_1, f_2$ are nonzero linear forms on $V$.
    Since $f_1, f_2 \ne 0$ there exist $x, y \in V$ such that
    $f_1(x) \ne 0$ and $f_2(y) \ne 0$.
    Let $v = x + y$. Then
    \[
        f(v, v) = f_1(v)f_2(v) \ne 0,
    \]
    which contradicts $f$ being alternate.
\end{proof}

\begin{proof}
    Suppose $n = 2$ and $f_1, f_2$ are linearly independent linear forms on $V$.
    If $\ker f_1 \subseteq \ker f_2$ then $f_2$ is a scalar multiple of $f_1$,
        contradicting linear independence.
    Thus there exists $x \in \ker f_1 \setminus \ker f_2$.
    Similarly there exists $y \in \ker f_2 \setminus \ker f_1$.
    Then
    \[f(x, y) = f_1(x)f_2(y) - f_1(y)f_2(x) = -f_1(y)f_2(x) \ne 0,\]
    so $f$ is not identically zero. For any $x \in V$,
    \[f(x, x) = f_1(x)f_2(x) - f_1(x)f_2(x) = 0,\]
    so $f$ is alternate.
\end{proof}

\begin{proof}
    Suppose $x_1, \ldots, x_n$ is a basis of $V$
        and $f_1, \ldots, f_n$ is the dual basis of $V'$.
    We show $f(x_1, \ldots, x_n) = 1$, after which Lemma 7.1.12 guarantees uniqueness.
    Since $f_i(x_j) = \delta_{ij}$, the only permutation $\sigma$ for which
    $f_{\sigma(1)}(x_1) \cdots f_{\sigma(n)}(x_n) \ne 0$ is $\sigma = \iota$,
    giving $(\operatorname{sgn} \iota) \cdot 1 \cdots 1 = 1$.
\end{proof}

\begin{tcolorbox}[title=Problem 3, breakable]
    
\end{tcolorbox}